% Section: properties-accessibility

\documentclass[../../../include/open-logic-section]{subfiles}

\begin{document}

\olfileid[fr]{nml}{frd}{acc}

\olsection{Propriétés des relations d'accessibilité}

De nombreuses !!{formula}s modales caractérisent des propriétés
simples, et parfois familières, de la relation d'accessibilité.
Dans un premier sens, tout modèle dont la relation possède une
propriété donnée rend vraie la !!{formula} correspondante,
ainsi que toutes ses instances par substitution. Commençons par
cinq exemples classiques de propriétés des relations
d'accessibilité et des formules dont elles garantissent la vérité.

\begin{thm}\ollabel{thm:soundschemas}
  Soit $\mModel{M} = \tuple{W, R, V}$ un modèle. Si $R$ possède
  l'une des propriétés indiquées à gauche dans le
  \olref{tab:five}, toute instance de la !!{formula} correspondante,
  indiquée à droite, est vraie dans~$\mModel{M}$.
\end{thm}

\begin{table}[t]
    \begin{tabular}{| l || l |}
      \hline
      {\emph{Si $R$ est \dots}} & {\emph{alors \dots est vraie dans~$\mModel{M}$ :}} \\
      \hline \hline
      \emph{sérielle} : $\forall u \exists v Ruv$ & \hbox to.43\textwidth
           {$\Box p \lif \Diamond p$ \hfill (\Ax{D})} \\
      \hline
      \emph{réflexive} : $\forall w Rww$
      & \hbox to.43\textwidth
          {$\Box p \lif p$ \hfill (\Ax{T})} \\
      \hline
      \emph{symétrique} : &  \hbox to .43\textwidth
           {$p \lif \Box\Diamond p$ \hfill (\Ax{B})} \\
           $\forall u\forall v(Ruv \lif Rvu)$ & \\
      \hline
      \emph{transitive} : & \hbox to.43\textwidth
           {$\Box p \lif \Box \Box p$ \hfill (\Ax{4})} \\
      $\forall u \forall v \forall w ((Ruv \land Rvw) \lif Ruw)$ & \\
      \hline 
      \emph{euclidienne} : & \hbox to .43\textwidth
        {$\Diamond p \lif \Box\Diamond p$ \hfill (\Ax{5})} \\
      $\forall w \forall u \forall v ((Rwu \land Rwv) \lif Ruv)$ &\\
      \hline
    \end{tabular}
    \caption{Cinq faits de correspondance.}
    \ollabel{tab:five}
  \end{table} 


\begin{proof}
  Traitons le cas de \Ax{B}. Pour établir qu'un schéma est vrai
  dans un modèle, il faut montrer que chacune de ses instances
  est vraie en tous les mondes de ce modèle. Soit donc
  $!A \lif \Box\Diamond !A$ une instance quelconque de \Ax{B},
  et soit $w \in W$ un monde arbitraire. Supposons que
  l'antécédent $!A$ soit vrai en $w$; montrons que
  $\Box \Diamond !A$ l'est aussi en $w$. Il suffit de montrer
  que $\Diamond !A$ est vraie en tout monde $w'$ accessible
  depuis $w$. Or, pour tout $w'$ tel que $Rww'$, la symétrie
  donne $Rw'w$ (voir la \olref{fig:Bsymm}). Puisque
  $\mSat{M}{!A}[w]$, nous avons
  $\mSat{M}{\Diamond !A}[w']$. Le monde $w'$ tel que $Rww'$
  étant arbitraire, $\mSat{M}{\Box\Diamond!A}[w]$.

  Nous laissons les autres cas en exercice.
\end{proof}

\begin{prob}
  Complétez la preuve du
  \olref[nml][frd][acc]{thm:soundschemas}.
\end{prob}

\begin{figure}
  \begin{center}
    \begin{tikzpicture}[modal]
      \node[world] (w1) [label={below:$\mSat{{}}{\formula{A}}$\\
          $\mSat{{}}{\Box\Diamond \formula{A}}$}] {$w$}; 
      \node[world] (w2) [label=below:{$\mSat{{}}{\Diamond \formula{A}}$},
        right=of w1] {$w'$}; 
      \draw[->, bend left] (w1) to (w2); 
      \draw[->, bend left] (w2) to (w1); 
    \end{tikzpicture}
  \end{center}
\caption{L'argument fondé sur la symétrie.}
\ollabel{fig:Bsymm}
\end{figure}

Les réciproques du \olref{thm:soundschemas} ne valent pas en général :
le fait qu'un modèle vérifie un schéma n'implique pas que sa
relation d'accessibilité possède la propriété correspondante.
Pour \Ax{T} et la réflexivité, considérons d'abord un modèle
où le schéma échoue : posons $W = \{w\}$,
$R = \emptyset$ et $V(p) = \emptyset$. La relation $R$
n'est pas réflexive; pourtant $\mSat{M}{\Box p}[w]$ et
$\mSat/{M}{p}[w]$, si bien que cette instance de \Ax{T}
est fausse dans~$\mModel{M}$. D'autres instances, par exemple
$\Box \lnot p \lif \lnot p$, y sont vraies. Il est moins
immédiat de construire un modèle non réflexif~$\mModel{M}$
où \emph{toutes les instances par substitution} de~\Ax{T}
sont vraies. De tels modèles existent néanmoins :

\begin{prop}\ollabel{prop:reflexive}
  Soit $\mModel{M} = \tuple{W, R, V}$ un modèle tel que
  $W = \{u, v\}$ et $R = \{(u,v),(v,u)\}$ : les deux mondes
  sont reliés l'un à l'autre, mais aucun à lui-même. Supposons
  en outre que, pour tout $p$,
  $u \in V(p) \Leftrightarrow v \in V(p)$. Alors :
  \begin{enumerate}
  \item Pour toute $!A$, $\mSat{M}{!A}[u]$ si et seulement si
    $\mSat{M}{!A}[v]$ (procédez par induction sur $!A$).
  \item Toute instance de~\Ax{T} est vraie dans $\mModel{M}$.
  \end{enumerate}
  Puisque $\mModel{M}$ n'est pas réflexif (sa relation est même
  \emph{irréflexive}), la réciproque du
  \olref{thm:soundschemas} échoue pour \Ax{T}. Des arguments
  analogues valent pour certains autres schémas du
  \olref{thm:soundschemas}, mais pas pour tous.
\end{prop}

\begin{prob}
  Démontrez les assertions de la
  \olref[nml][frd][acc]{prop:reflexive}.
\end{prob}

Nous nous concentrerons sur les cinq !!{formula}s classiques
\Ax{D}, \Ax{T}, \Ax{B}, \Ax{4} et \Ax{5}, mais
le \olref{tab:anotherfive} présente quelques autres propriétés des
relations d'accessibilité. La relation~$R$ est
\emph{partiellement fonctionnelle} si chaque monde donne accès
à au plus un monde. Si chaque monde donne accès à exactement
un monde, nous disons qu'elle est \emph{fonctionnelle}.
Les relations fonctionnelles sont donc précisément celles qui
sont à la fois sérielles et partiellement fonctionnelles.
Ce terme indique que la relation d'accessibilité se comporte
comme une fonction, éventuellement partielle. Une relation
est \emph{faiblement dense} si, chaque fois que $Ruv$, il existe
un monde~$w$ « entre » $u$ et~$v$. En un certain sens, la
densité faible inverse la transitivité : dans une relation
transitive, si l'on peut atteindre $v$ depuis~$u$ en passant
par~$w$, on peut aussi l'atteindre directement; dans une
relation faiblement dense, si l'on peut atteindre $v$
directement depuis~$u$, on peut également y parvenir en passant
par un certain monde~$w$. Une relation est
\emph{faiblement dirigée} si, chaque fois que l'on peut
atteindre les mondes $u$ et~$v$ depuis un monde~$w$, on peut
atteindre un même monde~$t$ depuis $u$ et depuis~$v$.
On parle aussi parfois de « propriété du diamant » ou de
« confluence ».

\begin{table}[t]
    \begin{tabular}{| p{.50\textwidth} || p{.43\textwidth} |}
      \hline
      {\emph{Si $R$ est \dots}} & {\emph{alors \dots est vraie dans~$\mModel{M}$ :}} \\
      \hline \hline
      \emph{partiellement fonctionnelle} : &
      \multirow{2}{*}{$\Diamond p \lif \Box p$} \\
      $ \forall w \forall u \forall v ((Rwu \land Rwv) \lif u =v )$ & \\
      \hline
      \multirow{2}{*}{\emph{fonctionnelle} : $\forall w \exists v \forall u(Rwu \liff \eq[u][v]) $}
      & \multirow{2}{*}{$\Diamond p \liff \Box p$} \\
      & \\
      \hline
      \emph{faiblement dense} : &  \multirow{2}{*}{$\Box \Box p \lif
        \Box p$} \\
      $\forall u\forall v(Ruv \lif \exists w(Ruw \land Rwv))$ & \\
      \hline
      \emph{faiblement connexe} : &
      \multirow{3}{*}{\hbox to.43\textwidth{$
        \begin{array}{@{}l@{}}
          \Box ((p \land \Box p) \lif q) \lor {} \\
          \qquad \Box ((q \land \Box q) \lif p)
        \end{array}$ \hfill (\Ax{L})}
      } \\
      $\forall w \forall u \forall v ((Rwu \land Rwv) \lif {}$ &\\
      \qquad $(Ruv \lor u=v \lor Rvu))$ & \\
      \hline 
      \emph{faiblement dirigée} : & \multirow{3}{*}{\hbox to .43\textwidth
        {$\Diamond\Box p \lif \Box\Diamond p$\hfill (\Ax{G})}} \\
      $\forall w \forall u \forall v ((Rwu \land Rwv) \lif {}$ & \\
      \qquad $\exists t (Rut \land Rvt))$ & \\
      \hline
    \end{tabular}
    \caption{Cinq autres faits de correspondance.}
    \ollabel{tab:anotherfive}
  \end{table} 

\begin{prob}
  Soit $\mModel{M} = \tuple{W, R, V}$ un modèle. Montrez que,
  si $R$ satisfait l'une des propriétés indiquées à gauche
  dans le \olref[nml][frd][acc]{tab:anotherfive}, toute instance
  de la formule correspondante, indiquée à droite, est vraie
  dans~$\mModel{M}$.
\end{prob}
\begin{editorial}Les relations des deux exemples non réflexifs sont précisées : dans le singleton, la relation est vide ; dans le modèle à deux mondes, elle contient exactement les deux flèches entre mondes distincts. La source ne précisait pas ces exclusions des boucles. Elles sont nécessaires pour que ces exemples aient bien la non-réflexivité annoncée.\end{editorial}


\end{document}
